\section{Claims, evidence and boundaries}
\label{sec:claim-scope}
Table~\ref{tab:claim-scope} separates the conclusions established by this artifact from the stronger claims that require a different experiment. A negative result has a named intervention and baseline; a positive result has a measured service objective and an explicit execution boundary.

\begin{table}[htbp]\centering\small
\caption[What each result establishes]{A guide to the interpretation of the current evidence. All timing claims in this table concern simulated DRAM service.}\label{tab:claim-scope}
\begin{tabular}{>{\raggedright\arraybackslash}p{0.26\textwidth}>{\raggedright\arraybackslash}p{0.34\textwidth}>{\raggedright\arraybackslash}p{0.29\textwidth}}
\toprule
Claim & Supporting comparison & Boundary\\\midrule
A row-hit improvement can accompany a timing loss. & Identical-work microbenchmarks, query-aware/BFS contrasts and command accounting. & Does not imply that row reuse is generally harmful.\\
Established placements can outperform the local heuristic. & RCM, Gorder and Porder under matching flat and bounded replay. & Does not reproduce their cache-aware application benchmarks.\\
Interleaving available vector chunks can reduce service. & Original/sorted/round-robin on matching expansions, including paired and per-query effects. & Depends on footprint, mapping, available work and omitted overhead.\\
A training profile can guide the fetch schedule. & Fixed 50-query rule, held-out comparisons to every fixed order, oracle gaps and calibration cost. & Exploratory, one corpus and operating point; future-query gains are not guaranteed.\\
A configurable near-memory gather path merits implementation testing. & Conditional service savings and an explicit buffering/execution sketch. & No implemented accelerator, end-to-end speedup or energy result.\\\bottomrule
\end{tabular}
\end{table}
